\section{Results and Discussion}\label{sec:synthesis}

The three preceding sections measure the same sixteen languages along
independent dimensions: the characters in which they are written, the sounds in
which they are spoken, and the terrain across which they are distributed.
Considered jointly, these measurements sustain an inference that none of them
supports in isolation. The present section develops that inference through a
systematic comparison of the two quantitative matrices, examines whether the
comparison survives the methodological objections to which it is exposed, and
concludes by identifying the findings upon which all three analyses converge.

\subsection{Comparative Structure of the Two Matrices}

A comparison of \cref{fig:ortho_sim_matrix,fig:phonetic_sim_matrix} indicates
that the orthographic and phonetic measures agree closely with respect to the
ordering of language relationships while differing substantially with respect to
their scale. Across the 120 distinct language pairs, the two measures rank those
pairs in a near-identical sequence, yielding a Spearman rank correlation of
\( \rho = 0.83 \) (\( p < 10^{-30} \)). The pairs that are orthographically
proximate are, with few exceptions, the same pairs that are phonetically
proximate, and the two measures may therefore be understood as recording a
common underlying ordering rather than two unrelated properties.

The distributions themselves, however, are not commensurable, as
\cref{tab:metric_comparison} sets out. Orthographic similarity is centred on
0.6596 and occupies a range of 0.5968, whereas phonetic similarity is centred on
0.8985 and occupies a range of 0.2305, approximately 39 per cent of the
orthographic figure. The relationships that both measures identify are
consequently registered in a considerably more attenuated form in speech than in
writing.

\begin{table}[h!]
    \centering
    \scriptsize
    \setlength{\tabcolsep}{3.5pt}
    \begin{tabular}{lrrrrr}
        \hline
                       &                & \multicolumn{2}{c}{\textbf{Orthographic}} & \multicolumn{2}{c}{\textbf{Phonetic}}                             \\
        \textbf{Pairs} & \textbf{\(n\)} & \textbf{Mean}                             & \textbf{Range}                        & \textbf{Mean} & \textbf{Range} \\
        \hline
        All            & 120            & 0.6596                                    & 0.5968                                & 0.8985        & 0.2305         \\
        Complete       & 28             & 0.6909                                    & 0.4473                                & 0.8956        & 0.2305         \\
        NT only        & 28             & 0.6262                                    & 0.5282                                & 0.9008        & 0.1495         \\
        Mixed          & 64             & 0.6605                                    & 0.5819                                & 0.8987        & 0.1976         \\
        \hline
    \end{tabular}
    \caption{Distribution of off-diagonal similarity under each measure, overall
        and stratified by the completeness of the two corpora compared.
        \emph{Complete} denotes pairs in which both corpora are complete Bibles,
        \emph{NT only} pairs in which both are New Testaments, and \emph{Mixed}
        one of each.}\label{tab:metric_comparison}
\end{table}

\subsection{Robustness of the Scale Difference}

Because the two measures are computed by different procedures, a Jaccard
coefficient over trigram sets in the orthographic case and a cosine over
frequency vectors in the phonetic, the possibility must be considered that the
difference in scale is an artefact of the instruments rather than a property of
the languages. The Jaccard coefficient is in principle sensitive to disparities
in set size, and eight of the sixteen corpora comprise only the New Testament
(\cref{tab:seeds}); corpus size therefore presents itself as a candidate
confound, and one that would, if operative, inflate the orthographic range
precisely as observed.

Neither the absolute difference in
corpus size between the members of a pair (\( \rho = 0.06 \) for the
orthographic measure, \( \rho = 0.03 \) for the phonetic) nor the size of the
smaller corpus in a pair (\( \rho = 0.09 \) and \( \rho = -0.01 \)
respectively) is appreciably associated with either similarity score. The
disparity in range moreover persists when the comparison is restricted to pairs
of equal corpus completeness, as the lower rows of \cref{tab:metric_comparison}
show: among the 28 pairs in which both corpora are complete Bibles the
orthographic range remains approximately twice the phonetic, and among the 28
pairs in which both are New Testaments it remains more than three times the
phonetic. The compression of the phonetic scale relative to the orthographic is
accordingly a property of what is being measured rather than of the corpora from
which the measurements derive.

\subsection{Interpretation}

The conjunction of close agreement in ordering with substantial disagreement in
scale is consistent with the two layers differing in age. Every language in the
sample acquired the Latin script within approximately the last four centuries,
through a colonial encounter that varied by region, by missionary order, and by
the period in which each translation was standardised. Orthography constitutes
the most recent and the most administratively tractable layer of a language, and
it registers those variations at full amplitude. Phonology is subject to neither
condition, having been established by no administrative decision and having
diverged from a common Austronesian ancestor over a period of millennia before
any of these languages was committed to writing. The narrower phonetic band may
therefore be understood as the residue of that shared inheritance, which remains
legible beneath four centuries of divergent orthographic practice. The colonial
imposition of writing did not reconfigure the relationships obtaining among
these languages; it amplified distinctions that phonology continues to preserve
in a more attenuated form.

\subsection{Convergent Findings}

Three findings recur across all three analyses and may be regarded as
comparatively secure, in that they rest upon measurements that could have
disagreed and did not.

Chavacano occupies a peripheral position under every measure. It is the most
orthographically distant language in the sample, with a mean similarity of
0.4190, and likewise the most phonetically distant, at 0.8414, while the
geographic analysis independently identifies Zamboanga as a contact zone in
which Spanish merged with Cebuano and Tausug substrates. Tausug furnishes the
comparison that lends this observation its force: comparably remote from the
remainder of the sample and comparably exposed to external contact, it
nonetheless remains phonetically central. Geographic distance and language
contact are therefore not of themselves sufficient to restructure a sound
system, whereas creolisation demonstrably is.

Romblomanon occupies a central position under every measure. It records the
highest mean orthographic similarity in the sample, at 0.7468, together with the
highest mean phonetic similarity, at 0.9174, and the geographic analysis places
Romblon between the Luzon and Visayan spheres. That a language situated at a
crossroads should resemble the varieties lying on either side of it, in writing
as in speech, is the expected consequence of sustained contact with both.

Northern Luzon constitutes a distinct zone under every measure. Ilocano is the
second most peripheral language both orthographically, at 0.5858, and
phonetically, at 0.8491, while remaining proximate to Paranan and Pangasinan,
and the Cordillera and Sierra Madre ranges appear in \cref{fig:lingmap_out} at
precisely the boundary that division implies. In this instance the three
analyses converge upon a single explanation, and that explanation is
topographic.

The centre--periphery ordering is in general preserved between the two measures
(\( \rho = 0.82 \) over per-language mean similarity), with Chavacano, Ilocano,
and Kapampangan the three most peripheral languages under both, and Kinaray-a,
Bikol, and Romblomanon the three most central under both.

\subsection{Systematic Divergences}

The disagreements between the measures are as instructive as the convergences,
and they are not randomly distributed. The pairs that are most similar
phonetically relative to their orthographic ranking are consistently pairs of
geographically remote languages: Kapampangan and Waray-Waray, separated by
Central Luzon and the Eastern Visayas, rank fourteenth of 120 orthographically
but seventy-sixth phonetically; Paranan and Tausug, separated by the length of
the archipelago, rank thirty-ninth and ninety-fourth respectively; and Yami and
Tausug, occupying its northern and southern extremities, rank thirtieth and
eighty-fourth. The Paranan--Pangasinan affinity exhibits the same asymmetry,
being strong phonetically at 0.9677 and merely moderate orthographically at
0.7828, and Kapampangan is phonetically closer to the distant Kinaray-a, at
0.9456, than to adjacent Ilocano, at 0.8483.

That this divergence should concentrate among distant pairs is consistent with
the interpretation advanced above. Orthographic practice would be expected to
diverge most freely between communities whose writing was standardised
independently of one another, by different missionary orders, translation
committees, and regional traditions of literacy, whereas phonological
inheritance persists irrespective of any such decision. On this reading distance
is a proxy for that independence rather than its cause: remote communities are
simply less likely to have had their spelling fixed by the same hands. Cases of
this kind nevertheless mark the limits of what a similarity matrix computed over
trigrams can settle. They are equally consistent with the retention of
conservative features lost elsewhere, with convergent innovation, and with
historical contact along routes that the terrain does not record.
Distinguishing among these possibilities requires evidence that the present
study does not contain, namely cognate sets, reconstructed proto-forms, and
documented migration histories. The further limitations upon the measures
employed here are set out in \cref{sec:limitations}.
